% Part: computability
% Chapter: computability-theory
% Section: introduction

\documentclass[../../../include/open-logic-section]{subfiles}

\begin{document}

\olfileid{cmp}{thy}{int}
\olsection{مقدمة}

يتناول فرع المنطق المعروف بـ\emph{نظرية قابلية الحساب} المسائل المتصلة
بقابلية الدوال والمجموعات للحساب، أو بقابليتها النسبية للحساب. ومن دلائل
تأثير كليني أن الموضوع كان يعرف من قبل بـ\emph{نظرية العودية}؛ ويشيع
اليوم استعمال الاسمين كليهما.

لنسم الدالة~$f\colon\Nat\pto\Nat$ \emph{جزئية قابلة للحساب} إذا أمكن
حسابها في نموذج من نماذج الحساب. وإذا كانت $f$ كلية قلنا ببساطة إن
$f$ \emph{قابلة للحساب}. وتسمى العلاقة~$R$ ذات الدالة المميزة القابلة
للحساب~$\Char{R}$ قابلة للحساب أيضًا. وإذا كانت $f$ و~$g$ دالتين
جزئيتين، كتبنا $f(x)\fdefined$ لنعني أن $f$ معرفة عند~$x$، أي إن
$x$ في مجال~$f$؛ وكتبنا $f(x)\fundefined$ لنعني العكس، أي إن $f$
غير معرفة عند~$x$. وسنستعمل $f(x)\simeq g(x)$ لنعني أن $f(x)$ و
$g(x)$ إما غير معرفتين معًا، وإما معرفتان ومتساويتان معًا.

يمكن بحث نظرية قابلية الحساب من غير الاضطرار إلى الإحالة إلى نموذج
حساب معين. وللقيام بذلك نبين وجود دالة جزئية شاملة قابلة للحساب
~$\fn{Un}(k,x)$. ويتيح لنا هذا تعداد الدوال الجزئية القابلة للحساب.
وسنعتمد الترميز~$\cfind{k}$ للدلالة على الدالة الجزئية الأحادية الحجة
القابلة للحساب ذات الفهرس $k$، المعرفة بـ
$\cfind{k}(x)\simeq\fn{Un}(k,x)$. (استعمل كليني $\{k\}$ لهذا الغرض،
لكن هذا الترميز قل استعماله في الآونة الأخيرة.) وعلى نحو أعم قليلًا،
يمكننا أن نعدد بانتظام الدوال الجزئية القابلة للحساب ذات الرتب
المختلفة، وسنستعمل $\cfind{k}[n]$ للدلالة على الدالة العودية الجزئية
ذات $n$ حجة وذات الفهرس $k$ في التعداد.

إذا كانت $f(\vec x,y)$ دالة كلية أو جزئية، فإن
$\umin{y}{f(\vec x,y)}$ هي الدالة في~$\vec x$ التي تعيد أصغر~$y$
بحيث $f(\vec x,y)=0$، على افتراض أن $f(\vec x,0)$، …،
$f(\vec x,y-1)$ كلها معرفة؛ وإذا لم يوجد $y$ كهذا، كانت
$\umin{y}{f(\vec x,y)}$ غير معرفة. وإذا كانت $R(\vec x,y)$ علاقة،
عرفنا $\umin{y}{R(\vec x,y)}$ بأنه أصغر~$y$ بحيث تصدق
$R(\vec x,y)$؛ أي أصغر~$y$ بحيث $1\tsub\Char{R}(\vec x,y)=0$.

ولبيان أن دالة قابلة للحساب، ثمة طريقتان:
\begin{enumerate}
\item على نحو صارم: نصف آلة تورنغ أو دالة عودية جزئية وصفًا صريحًا،
  ونبين أنها تحسب الدالة المقصودة؛
\item على نحو غير صوري: نصف خوارزمية تحسبها، ونحتكم إلى أطروحة تشرش.
\end{enumerate}
ولا يوجد حد فاصل دقيق بين الطريقتين؛ فوصف مفصل لخوارزمية ينبغي أن يوفر
من المعلومات ما يكفي ليتضح نسبيًا كيف يمكن، من حيث المبدأ، تصميم آلة
تورنغ المناسبة أو متتالية التعريفات العودية الجزئية المناسبة. ولا يرجح
أن تكون التعريفات الصارمة تمامًا مفيدة، وسنحاول إيجاد حل وسط ملائم بين
النهجين؛ وباختصار، سنحاول إيجاد براهين حدسية ومع ذلك صارمة، يمكن الحصول
منها على التعريفات الدقيقة.

\end{document}
